\documentclass[a4paper]{article}

% 文章排版
\usepackage{indentfirst}
\setlength{\parindent}{0em}
\usepackage{autobreak}
\allowdisplaybreaks
\usepackage{parskip}
\usepackage{geometry}
\geometry{top=3cm,bottom=3cm,left=2cm,right=2cm}
\usepackage{anyfontsize}

% 数学物理宏包
\usepackage{amsmath}
\usepackage{physics}
\renewcommand\div\divisionsymbol
\DeclareMathOperator{\diag}{diag}
\DeclareMathOperator{\sgn}{sgn}
\newcommand{\trans}{\text{\textsf{T}}}
\usepackage{tabularray}
\UseTblrLibrary{amsmath}

\usepackage{xcolor}
\usepackage{fontspec}
\setmainfont{STIX Two Text}
\usepackage{unicode-math}
\unimathsetup{mathrm=sym,mathbf=sym,mathsf=sym,bold-style=ISO}
\setmathfont{STIX Two Math}

\usepackage[fontset=none]{ctex}
\setCJKmainfont{FZShuSong-Z01}[BoldFont=Source Han Sans CN,ItalicFont=FZKai-Z03]
\setCJKsansfont{Source Han Sans CN}

% 符号重定义
\AtBeginDocument{
    \let\uglyepsilon\epsilon
    \renewcommand\epsilon\varepsilon
    \let\uglyleq\leq
    \renewcommand\leq\leqslant
    \let\uglygeq\geq
    \renewcommand\geq\geqslant
}


% 题目
\title{\textbf{广义相对论（天文方向）\ \ \ \ 作业四}}
\author{国家天文台 张植竣}
\date{2025年10月13日}

\begin{document}

\maketitle

\begin{enumerate}
    \item 采用Euclid直角坐标系，伪球面的参数方程为：
    \[
        \begin{+cases}
            x = a \cos\theta \sech{u} \\
            y = a \sin\theta \sech{u} \\
            z = a \qty(u - \tanh{u})
        \end{+cases},\quad u \in (-\infty, +\infty),\ \theta \in [0, 2\pi)
    \]
    证明伪球面上的线元为：
    \[
        \dd{s}^2 = a^2 \qty(\tanh^2{u} \dd{u}^2 + \sech^2{u} \dd{\theta}^2)
    \]

    \textcolor{blue}{
    伪球面上的线元为：
    \[
        \dd{s}^2 = \dd{x}^2 + \dd{y}^2 + \dd{z}^2 = \qty(\pdv{x}{\theta} \dd{\theta} + \pdv{x}{u} \dd{u})^2 + \qty(\pdv{y}{\theta} \dd{\theta} + \pdv{y}{u} \dd{u})^2 + \qty(\pdv{z}{u} \dd{u})^2
    \]
    计算各偏导数：（注意$\tanh{u}$的导数为$\sech^2{u}$，$\sech{u}$的导数为$-\sech{u} \tanh{u}$）
    \[
        \pdv{x}{\theta} = -a \sin\theta \sech{u},\quad \pdv{x}{u} = -a \cos\theta \sech{u} \tanh{u}
    \]
    \[
        \pdv{y}{\theta} = a \cos\theta \sech{u},\quad \pdv{y}{u} = -a \sin\theta \sech{u} \tanh{u}
    \]
    \[
        \pdv{z}{u} = a \qty(1 - \sech^2{u}) = a \tanh^2{u}
    \]
    于是有：
    \[
        \dd{x} = -a \sin\theta \sech{u} \dd{\theta} - a \cos\theta \sech{u} \tanh{u} \dd{u}
    \]
    \[
        \dd{y} = a \cos\theta \sech{u} \dd{\theta} - a \sin\theta \sech{u} \tanh{u} \dd{u}
    \]
    \[
        \dd{z} = a \tanh^2{u} \dd{u}
    \]
    将各导数代入线元表达式，并利用三角恒等式$\sin^2\theta + \cos^2\theta = 1$，得到：
    \[
        \dd{s}^2 = a^2 \sech^2{u} \dd{\theta}^2 + a^2 \qty(\sech^2{u} \tanh^2{u} + \tanh^4{u}) \dd{u}^2
    \]
    再利用双曲函数恒等式$\sech^2{u} + \tanh^2{u} = 1$，可化简为：
    \[
        \dd{s}^2 = a^2 \qty(\tanh^2{u} \dd{u}^2 + \sech^2{u} \dd{\theta}^2)
    \]
    }

    \item 证明通过坐标变换，伪球面上的线元可以写成：
    \[
        \dd{s}^2 = a^2 \qty(\dd{X}^2 + \mathrm{e}^{-2|X|} \dd{\theta}^2)
    \]

    \textcolor{blue}{
    由上一问的结果：
    \[
        \dd{s}^2 = a^2 \qty(\tanh^2{u} \dd{u}^2 + \sech^2{u} \dd{\theta}^2) = a^2 \qty(\dd{X}^2 + \mathrm{e}^{-2|X|} \dd{\theta}^2)
    \]
    对比得：
    \[
        \tanh^2{u} \dd{u}^2 = \dd{X}^2,\quad \sech^2{u} = \mathrm{e}^{-2|X|}
    \]
    对第一式开方，取正号（因为$\tanh{u} > 0$时，$\dd{u}$与$\dd{X}$同号）：
    \[
        \tanh{u} \dd{u} = \dd{X}
    \]
    积分得：
    \[
        X = \ln(\cosh{u}) + C
    \]
    取$u = 0$时$X = 0$，则$C = 0$，所以：
    \[
        X = \ln(\cosh{u})
    \]
    或由第二式开方，取正号（因为$\sech{u} > 0$）：
    \[
        \sech{u} = \mathrm{e}^{-|X|} \implies \cosh{u} = \mathrm{e}^{|X|} \implies |X| = \ln(\cosh{u})
    \]
    取绝对值是因为这样可以将$X$定义在$(-\infty, +\infty)$上，而$\cosh{u} \geq 1$，等式右侧恒不小于0。 \\
    综上，令$|X| = \ln(\cosh{u})$，则：
    \[
        \dd{s}^2 = a^2 \qty(\dd{X}^2 + \mathrm{e}^{-2|X|} \dd{\theta}^2)
    \]
    }

    \item 基于问题2的线元，计算伪球面的Christoffel符号。

    \textcolor{blue}{
    由问题2的线元：
    \[
        \dd{s}^2 = a^2 \qty(\dd{X}^2 + \mathrm{e}^{-2|X|} \dd{\theta}^2)
    \]
    可得度规张量及其逆：
    \[
        g_{\mu\nu} = a^2 \mqty[1 & 0 \\ 0 & \mathrm{e}^{-2|X|}],\quad g^{\mu\nu} = \frac{1}{a^2} \mqty[1 & 0 \\ 0 & \mathrm{e}^{2|X|}]
    \]
    Christoffel符号的定义为：
    \[
        \Gamma^\sigma_{\mu\nu} = \frac{1}{2} g^{\sigma\rho} \qty(\partial_\mu g_{\nu\rho} + \partial_\nu g_{\mu\rho} - \partial_\rho g_{\mu\nu})
    \]
    由于度规张量为对角矩阵，Christoffel符号的计算可以简化为：
    \[
        \Gamma^0_{00} = \frac{1}{2} g^{00} \partial_0 g_{00},\quad \Gamma^0_{01} = \Gamma^0_{10} = \frac{1}{2} g^{00} \partial_1 g_{00},\quad \Gamma^0_{11} = - \frac{1}{2} g^{00} \partial_0 g_{11}
    \]
    \[
        \Gamma^1_{00} = -\frac{1}{2} g^{11} \partial_1 g_{00},\quad \Gamma^1_{01} = \Gamma^1_{10} = \frac{1}{2} g^{11} \partial_0 g_{11},\quad \Gamma^1_{11} = \frac{1}{2} g^{11} \partial_1 g_{11}
    \]
    这些分量原则上都是非零的，但由于$\partial_0 g_{00} = \partial_1 g_{00} = \partial_1 g_{11} = 0$，所以非零的分量只有：
    \[
        \Gamma^0_{11} = -\frac{1}{2} g^{00} \partial_0 g_{11} = -\frac{1}{2a^2} \cdot \partial_X \qty(a^2 \mathrm{e}^{-2|X|}) = -\frac{1}{2a^2} \qty(-2a^2 \mathrm{e}^{-2|X|} \sgn{X}) = \mathrm{e}^{-2|X|} \sgn{X}
    \]
    \[
        \Gamma^1_{01} = \Gamma^1_{10} = \frac{1}{2} g^{11} \partial_0 g_{11} = \frac{1}{2a^2} \mathrm{e}^{2|X|} \cdot \partial_X \qty(a^2 \mathrm{e}^{-2|X|}) = \frac{1}{2a^2} \mathrm{e}^{2|X|} \cdot \qty(-2a^2 \mathrm{e}^{-2|X|} \sgn{X}) = -\sgn{X}
    \]
    其中$\sgn{X}$为符号函数（绝对值函数的导数），定义为：
    \[
        \sgn{X} =
        \begin{cases}
            1, & X > 0 \\
            0, & X = 0 \\
            -1, & X < 0
        \end{cases}
    \]
    对于对角度规$g_{\mu\nu}$，Christoffel符号$\Gamma^\sigma_{\mu\nu}$只有在$\sigma$、$\mu$、$\nu$至少有两个值相同的情况下才可能非零，而对于二维空间，任意三个指标组合一定有两个是相同的。
    }
\end{enumerate}

\end{document}